\section{Runtime}
\label{sec:runtime}

The certificate of \S\ref{sec:theory} is conditional on the execution rules
E1--E4. The job of the runtime is to make those hypotheses true. It sits
between the agents and the target system as a middleware: every tool call an
agent emits is dispatched through it, and nothing reaches the target system
any other way. Figure~\ref{fig:arch} shows the path of a call.

\begin{figure}[t]
\centering
\begin{tikzpicture}[
  box/.style={draw=ccmuted, rounded corners=2pt, font=\scriptsize,
              minimum height=6mm, inner sep=3pt, align=center},
  arr/.style={-{Stealth[length=1.6mm]}, ccmuted, line width=0.6pt},
  lab/.style={font=\scriptsize, ccmuted}]
\node[box] (agent) at (0,0) {agent\\(model loop)};
\node[box, right=7mm of agent] (piece) {map call\\to piece};
\node[box, right=5mm of piece] (enf) {footprint\\check (E3)};
\node[box, right=5mm of enf] (lock) {piece\\locks (E1)};
\node[box, right=5mm of lock] (store) {target\\system};
\node[box, below=7mm of enf] (adm) {admission\\(Alg.~\ref{alg:admit})};
\draw[arr] (agent) -- node[above, lab] {call} (piece);
\draw[arr] (piece) -- (enf);
\draw[arr] (enf) -- (lock);
\draw[arr] (lock) -- (store);
\draw[arr] (enf) -- node[right, lab] {outside} (adm);
\draw[arr] (adm.west) -| node[pos=0.25, below, lab] {refuse / admit} (agent.south);
\end{tikzpicture}
\caption{The path of one tool call. A call inside the footprint of the
executing piece proceeds to locking. A call outside it is a plan revision: it
is admitted, possibly after coarsening, or refused and reported back to the
agent as a tool result.}
\label{fig:arch}
\end{figure}

\subsection{Mediation and fail-closed enforcement}

Every read, range scan and write is checked against the footprint of the
piece that is currently executing for that workflow. An access inside the
footprint proceeds. An access outside it never reaches the target system.
Instead it is treated as what it is, a proposal to revise the plan, and is
handed to Algorithm~\ref{alg:admit}. If the revision is admitted the plan is
updated and the call proceeds; if it is refused the call fails, and the agent
receives the refusal as the result of its tool call.

Enforcement does not depend on the agent's cooperation. This matters because
a model-driven agent does not follow a script: it can call any tool with any
argument at any time, and its own account of what it intends to do is not
evidence. The trust boundary is explicit. The tools and the middleware are
trusted; the agent that drives them is not.

\paragraph{Pieces for a model-driven agent.}
A scripted workflow declares its steps and their grouping directly. A
model-driven agent declares only a footprint per piece, and the runtime must
decide which piece a call belongs to. It advances from the piece currently
held to the first later piece whose footprint permits the access; if none
does, the access stays with the current piece and becomes a revision proposal.
Every gap between two tool calls of a model-driven agent is a reasoning gap,
so a piece that certification left unmerged closes when its call returns. A
piece that certification merged across a reasoning gap is held through it;
releasing it per call would undo exactly the repair that made the plan safe.

\subsection{Piece isolation: three locking disciplines}
\label{sec:runtime:locking}

Rule E1 requires each piece to be serializable with respect to every other
piece. We enforce it by locking, with read and write modes and with range
locks so that an insert into a scanned range is blocked. Locks never outlive a
piece: they are released when the piece commits, which is what lets a chopped
workflow give up its shared footprint before the model starts reasoning.

Within that frame there is a design choice---when a piece acquires its
locks---and we study three answers (Table~\ref{tab:locking}).

\begin{table}[t]
\caption{Three ways to isolate a piece. All release at the piece boundary and
all run under the same certificate.}
\label{tab:locking}
\small
\begin{tabular}{@{}l>{\raggedright}p{2.05cm}>{\raggedright}p{1.85cm}>{\raggedright\arraybackslash}p{1.9cm}@{}}
\toprule
 & acquires & on conflict & aborts when \\
\midrule
\cc     & whole footprint, one wait & wait & never \\
\cctwo  & per access, at first use  & wait & a wait-for cycle is detected \\
\ccs    & per access, at first use  & higher $\sigma$ waits; lower $\sigma$ wounds & a lower-$\sigma$ piece needs the lock \\
\bottomrule
\end{tabular}
\end{table}

\paragraph{\cc: atomic footprint.}
A piece acquires locks for its entire declared footprint atomically, in one
wait, before executing. Because nothing is acquired while holding a lock, no
wait-for cycle can form and the lock manager is deadlock-free by
construction. The cost is over-reservation in time: an object the piece will
not touch until its last step is locked from its first step.

\paragraph{\cctwo: two-phase locking inside a piece.}
A piece takes each lock when the access happens and releases all of them
together at the piece boundary: strict two-phase locking scoped to a piece.
Holds are shorter, but hold-and-wait returns, and with it deadlock. The
typical case is a lock upgrade. Two pieces each hold a read lock on an object
and each then requests the write lock. \cctwo detects the cycle in the
wait-for graph, aborts the member with the highest $\sigma$, undoes its writes
through the registered inverses, and restarts it.

\paragraph{\ccs: priority-ordered locking.}
The deadlock in \cctwo is the only thing standing between incremental locking
and a strict improvement, and it does not have to be detected. It can be made
unreachable. The device is the one \mtpo uses for notifications: a rank
$\sigma$ assigned to every workflow at launch, and a rule that influence flows
in one direction along it. Applied to lock waits, the rule is:
\begin{itemize}
\item if the requester's $\sigma$ is \emph{higher} than the holder's, the
  requester \emph{waits};
\item if the requester's $\sigma$ is \emph{lower} than the holder's, the
  holder is \emph{wounded}: it is aborted, its writes are undone, its locks
  are released, and the requester proceeds.
\end{itemize}
This is the wound-wait scheme of Rosenkrantz et al.~\cite{rosenkrantz1978}
with $\sigma$ as the timestamp. We chose it over wait-die because wait-die
aborts a requester the moment it meets any held lock, including one about to
be released.

The semantics line up with \mtpo as well. Under \mtpo the lower-ranked agent
writes as though it ran first and the higher-ranked agent absorbs the change
through a notification. Under \ccs the lower-ranked agent takes the lock as
though it ran first and the higher-ranked agent absorbs the change by aborting
and redoing. The order is the same and the direction of yielding is the same;
what differs is the mechanism of the yield---premise repair there, redo here.
\S\ref{sec:eval:agent} shows that this difference matters for agents.

\begin{proposition}[The certificate survives incremental locking]
\label{prop:atomic}
Under \cctwo and \ccs, every execution satisfies E1.
\end{proposition}

\begin{proof}
Both disciplines acquire a lock before each access and release no lock before
the piece commits or aborts, and an aborted piece's writes are undone before
its locks are released. This is strict two-phase locking with pieces as the
transactions, so the execution is conflict-serializable at piece granularity
and no piece observes an uncommitted write of another.
\end{proof}

One-shot acquisition was therefore never buying correctness. It was buying
deadlock freedom, and \ccs buys that differently:

\begin{proposition}[\ccs is deadlock-free]
\label{prop:acyclic}
Under \ccs the wait-for graph is acyclic at all times.
\end{proposition}

\begin{proof}
A workflow waits only for a workflow of strictly lower $\sigma$. A cycle in
the wait-for graph would require $\sigma$ to decrease strictly around the
cycle, which is impossible.
\end{proof}

The workflow with the lowest $\sigma$ never waits and is never wounded, so
some workflow always makes progress. A higher-ranked workflow can be wounded
repeatedly; the runtime caps restarts and reports a workflow that exhausts the
cap as failed. We measure how often wounds occur and whether the cap is ever
reached in \S\ref{sec:eval:locking}.

\subsection{Refusals, aborts and the agent}

Two runtime events are visible to the agent, and they ask different things of
it.

A \emph{refusal} is returned as the result of the offending tool call. The
agent's conversation is intact, and it is expected to continue: by retrying
once the blocking piece has released, by choosing an operation inside its
footprint, or by reporting that the intent cannot be met.

An \emph{abort}---a deadlock victim under \cctwo or a wound under
\ccs---restarts the agent. Its writes are undone, its locks are released, and
its conversation is reset to the original task statement. The agent does not
see a partial failure and is not asked to reason about one; it performs its
task again against the current state.

\subsection{Durable continuation}

Chopping trades failure atomicity for concurrency: once a piece commits, its
effects are visible even if the process dies before the next piece. The
certificate says nothing about this, and the runtime must. Each piece's writes
and the workflow's progress record, including its local variables, are
committed in one transaction of the underlying store. A crash inside that
transaction rolls back both, so pieces are atomic; a crash after it leaves a
durable marker from which recovery replays the remaining pieces forward (rule
E4).

\subsection{Implementations}
\label{sec:runtime:impl}

We built the mechanism twice. The \emph{durable runtime} is about 1{,}470
lines of Python over SQLite in write-ahead-logging mode. It runs workflows on
real threads, implements \cc with durable continuation, and is used for the
mechanism-level throughput, crash and certification-cost experiments.

The \emph{agent runtime} is a middleware in the style of
CoAgent~\cite{lyu2026coagent}: tools declare footprints and register
inverses, and a concurrency-control protocol is consulted on every call. In
it we implemented \cc, \cctwo and \ccs next to five other protocols---serial
execution, no coordination, whole-workflow 2PL and OCC with saga undo, and
\mtpo---so that all eight face the same agents, tools and workloads. Its
store is in memory, so it has no crash boundary and does not implement durable
continuation. An abort in this runtime restarts the whole workflow and
compensates the writes of pieces it had already committed, rather than
restarting only the aborted piece from a progress record; \S\ref{sec:discussion}
discusses the consequence.
